%% Question bank chapter file - generated by bank_tools.py build (do not edit while the bank is still being built:
%% it is regenerated from the per-test data in Question Bank\data). Every question is one bankq block:
%%   \begin{bankq}[id=..., source={...}, paper=.., qno=.., type=..., subtopic={...}, answer={...}, answer_src=key|solved,
%%                 working={...}, difficulty=1-10, diff_parts={K S I C F}, diff_note={...},
%%                 group=<passage id or empty>, also_in={duplicate copies in other tests}]
%% difficulty = min(10, K+S+I+C+F-1): K concepts 1-3, S steps 1-4, I insight 0-2, C computation 0-2, F format 0-1
%%   ...question LaTeX, options, figures (\bfig / \bfigi -> img/)...
%%   \end{bankq}
%% Passages shared by several questions are bankpassage blocks; their questions carry group=<passage id>.
\documentclass[10pt]{article}
\usepackage[a4paper,margin=1.8cm]{geometry}
\usepackage[T1]{fontenc}
\usepackage{amsmath,amssymb,mathtools}
\usepackage{graphicx}
\usepackage[export]{adjustbox}
\usepackage{xcolor,needspace,keyval}
\usepackage{enumitem}
\usepackage{booktabs,array,tabularx}
\usepackage[hidelinks]{hyperref}
\setlength{\parindent}{0pt}
\setlength{\parskip}{3pt}
\makeatletter
\define@key{bq}{id}{\def\bq@id{#1}}\define@key{bq}{source}{\def\bq@src{#1}}\define@key{bq}{paper}{\def\bq@paper{#1}}
\define@key{bq}{qno}{\def\bq@qno{#1}}\define@key{bq}{type}{\def\bq@type{#1}}\define@key{bq}{subtopic}{\def\bq@sub{#1}}
\define@key{bq}{answer}{\def\bq@ans{#1}}\define@key{bq}{answer_src}{\def\bq@asrc{#1}}\define@key{bq}{working}{\def\bq@work{#1}}
\define@key{bq}{group}{\def\bq@group{#1}}\define@key{bq}{also_in}{\def\bq@also{#1}}
\define@key{bq}{difficulty}{\def\bq@diff{#1}}\define@key{bq}{diff_parts}{\def\bq@dparts{#1}}\define@key{bq}{diff_note}{\def\bq@dnote{#1}}
\newenvironment{bankq}[1][]{\def\bq@work{}\def\bq@also{}\def\bq@group{}\def\bq@diff{?}\def\bq@dparts{}\def\bq@dnote{}\setkeys{bq}{#1}%
  \par\needspace{4\baselineskip}\medskip\noindent{\small\textbf{\bq@id}\hfill\color{gray}[\bq@src, Paper \bq@paper, Q\bq@qno]
  \bq@type\ $\cdot$ \bq@sub\ $\cdot$ \textbf{D\bq@diff}}\par\nopagebreak\relax}
  {\par{\small\color{teal}Answer: \textbf{\bq@ans}\ifx\bq@work\@empty\else\ (solved: \bq@work)\fi
  \ifx\bq@also\@empty\else\par\color{gray}Also in: \bq@also\fi}\par}
\define@key{bp}{id}{\def\bp@id{#1}}\define@key{bp}{source}{\def\bp@src{#1}}\define@key{bp}{questions}{\def\bp@qs{#1}}
\newenvironment{bankpassage}[1][]{\setkeys{bp}{#1}\par\needspace{4\baselineskip}\bigskip\noindent
  {\small\textbf{Passage \bp@id}\hfill\color{gray}[\bp@src; questions \bp@qs]}\par\nopagebreak\relax\itshape}{\par}
\makeatother
\newcommand{\bfig}[2]{\par\smallskip{\centering\includegraphics[width=#1pt,max width=0.62\linewidth,max height=0.32\textheight]{img/#2}\par}\smallskip}
\newcommand{\bfigi}[2]{\includegraphics[width=#1pt,max width=\linewidth,max height=0.32\textheight,valign=c]{img/#2}}
\newcommand{\bmissing}[1]{\textit{[figure missing in the source conversion: \detokenize{#1}]}}
\newenvironment{qoptions}{\begin{itemize}[leftmargin=1.8em,itemsep=1pt,topsep=2pt,label={}]}{\end{itemize}}
\newcommand{\opt}[2]{\item[(#1)] #2}
\providecommand{\fwstart}{}\providecommand{\fwend}{}\providecommand{\qfullsplit}{}
\newcolumntype{L}[1]{>{\raggedright\arraybackslash}p{#1}}
\newcolumntype{C}[1]{>{\centering\arraybackslash}p{#1}}
\newcolumntype{R}[1]{>{\raggedleft\arraybackslash}p{#1}}
\providecommand{\sectionhdr}[1]{}\providecommand{\subsectionhdr}[1]{}


\begin{document}
{\Large\bfseries Indefinite Integration --- Question Bank}\par\smallskip
{\small Mathematics chapter M26. Questions from all converted test papers, grouped by source series. Answers: printed key where it matches the paper, otherwise solved (answer\_src=solved, with working). Exact repeats across papers appear once, with also\_in listing the other copies.}\par\medskip\hrule

\section*{TOAS papers}
\begin{bankq}[id=TOAS09-P1-MATH-S1-Q01, source={TOAS - Test 09 (30 May 2022)}, paper=1, qno=1, type=single, subtopic={By parts}, answer={D}, answer_src=key, difficulty=3, diff_parts={K1 S2 I1 C0 F0}, diff_note={concepts: integration by recognising a derivative | steps: simplify e\^{}(ln x) factor; guess x\^{}2 e\^{}(tan\^{}-1 x + x\^{}2); verify by differentiation | I: integrand is the derivative of a product}]
{\small\[\begin{aligned}\int\Big(&2e^{\tan^{-1}x+x^2+\ln x}\\&+\frac{2x^3+2x+1}{x+\frac1x}\,e^{\tan^{-1}x+x^2+\ln x}\Big)dx\end{aligned}\]}
is equal to (where C is constant of integration)
\begin{qoptions}
\opt{A}{$\dfrac{x^3}{3}e^{\tan^{-1}x+x^2+\ln x}+C$}
\opt{B}{$\dfrac{x^2}{2}e^{\tan^{-1}x+x^2+\ln x}+C$}
\opt{C}{$xe^{\tan^{-1}x+x^3+\ln x}+C$}
\opt{D}{$x^2e^{\tan^{-1}x+x^2}+C$}
\end{qoptions}
\end{bankq}
\begin{bankq}[id=TOAS09-P2-MATH-S1-Q04, source={TOAS - Test 09 (30 May 2022)}, paper=2, qno=4, type=single, subtopic={Partial fractions}, answer={A}, answer_src=key, difficulty=3, diff_parts={K1 S2 I1 C0 F0}, diff_note={concepts: partial fractions integration | steps: split into A/x and a term in x\^{}2; integrate; combine logs | I: write the numerator with the derivative of the denominator}]
$\displaystyle\int\dfrac{2\left(x^3-1\right)}{x\left(2x^3+1\right)}dx$ is equal to- (where C is the constant of integration)
\begin{qoptions}
\opt{A}{$\ell n\left(2x+\dfrac1{x^2}\right)+C$}
\opt{B}{$\ell n\left(2x^3+1\right)+C$}
\opt{C}{$\tan^{-1}\dfrac{1}{x^{3/2}+1}+\ell n\left(\dfrac{x-1}{x+1}\right)+C$}
\opt{D}{$\ell n\left(x+\sqrt{\dfrac{x-1}{x+1}}\right)+C$}
\end{qoptions}
\end{bankq}
\begin{bankq}[id=TOAS09-P2-MATH-S1-Q09, source={TOAS - Test 09 (30 May 2022)}, paper=2, qno=9, type=single, subtopic={By parts}, answer={A}, answer_src=key, difficulty=2, diff_parts={K1 S1 I1 C0 F0}, diff_note={concepts: integration of trigonometric products | steps: integrate by parts or match coefficients with f, g linear | I: undetermined coefficients}]
If linear functions $f(x)$ and $g(x)$ satisfy $\displaystyle\int\left[(1-2x)\cos x+(3+2x)\sin x\right]dx=f(x)\sin x+g(x)\cos x+C$ (where C is constant of integration), then
\begin{qoptions}
\opt{A}{$f(x)=3-2x$}
\opt{B}{$g(x)=5-2x$}
\opt{C}{$f(x)=3+2x$}
\opt{D}{$g(x)=-5+2x$}
\end{qoptions}
\end{bankq}
\begin{bankq}[id=TOAS09-P2-MATH-S1-Q14, source={TOAS - Test 09 (30 May 2022)}, paper=2, qno=14, type=single, subtopic={Substitution}, answer={A}, answer_src=key, difficulty=4, diff_parts={K1 S2 I2 C0 F0}, diff_note={concepts: integration by recognising a derivative | steps: compare with the derivative of x\^{}28 (1 + x + x\^{}2)\^{}7/7; verify; select the option | I: recognise the derivative of a product}]
$\displaystyle\int\left(6x^2+5x+4\right)\left(x^2+x+1\right)^6\cdot x^{27}dx$ equals (where C is integration constant)
\begin{qoptions}
\opt{A}{$\dfrac{x^{28}\left(1+x+x^2\right)^7}{7}+C$}
\opt{B}{$\dfrac{x^{24}\left(1+x+x^2\right)^7}{7}+C$}
\opt{C}{$\dfrac{x^{24}\left(1+x+x^2\right)^6}{7}+C$}
\opt{D}{$\dfrac{x^{28}\left(1+x+x^2\right)^8}{7}+C$}
\end{qoptions}
\end{bankq}
\begin{bankq}[id=TOAS09-P2-MATH-S2-Q10, source={TOAS - Test 09 (30 May 2022)}, paper=2, qno=10, type=numerical, subtopic={Substitution}, answer={2}, answer_src=solved, working={Integrand $=e^{4x}/(e^{4x}+1)^{1/4}$; $u=e^{4x}+1$ gives $\frac13(e^{4x}+1)^{3/4}$: $k=3,\alpha=4,\beta=3,\gamma=4$; $1+1=2$}, difficulty=4, diff_parts={K1 S2 I1 C0 F1}, diff_note={concepts: integration by substitution | steps: simplify to e\^{}(4x)/(e\^{}(4x) + 1)\^{}(1/4); substitute u = e\^{}(4x) + 1; integrate | I: combine exponentials}]
If $\displaystyle\int\dfrac{\sqrt{e^{7x}}}{\sqrt[4]{e^{2x}+e^{-2x}}}dx=\dfrac1k\left(e^{\alpha x}+1\right)^{\beta/\gamma}+C$ (where C is constant of integration \& $k$, $\alpha$, $\beta$, $\gamma$ are positive integers and $\beta$, $\gamma$ are relatively prime) then $\left(\dfrac k\beta+\dfrac\gamma\alpha\right)$ is
\end{bankq}
\begin{bankq}[id=TOAS10-P1-MATH-S1-Q02, source={TOAS - Test 10 (12 Jun 2022)}, paper=1, qno=2, type=multi, subtopic={Substitution}, answer={B}, answer_src=key, difficulty=3, diff_parts={K1 S1 I1 C0 F1}, diff_note={concepts: integral from a derivative form | steps: differentiate sqrt(x\^{}4 - x\^{}3 + 2)/x\^{}2; compare with the integrand to find k | I: test the given form by differentiation}]
If $\displaystyle\int\dfrac{x^3-8}{x^3\sqrt{x^4-x^3+2}}\,dx=\dfrac{k\sqrt{x^4-x^3+2}}{x^2}+C$, then the value of $k$ is
\begin{qoptions}
\opt{A}{3}
\opt{B}{2}
\opt{C}{7}
\opt{D}{4}
\end{qoptions}
\end{bankq}
\begin{bankq}[id=TOAS10-P1-MATH-S1-Q03, source={TOAS - Test 10 (12 Jun 2022)}, paper=1, qno=3, type=multi, subtopic={Standard forms}, answer={B,C,D}, answer_src=key, difficulty=5, diff_parts={K1 S3 I1 C0 F1}, diff_note={concepts: integration of a rational function over a radical | steps: differentiate the right side; multiply by the radical and expand; compare coefficients to find a; find b; check the options | I: differentiate the claimed antiderivative}]
If $\displaystyle\int\dfrac{3z^3-8z+5}{\sqrt{z^2-4z-7}}\,dz=(z^2+az+36)\sqrt{z^2-4z-7}+b\ln\left|z-2+\sqrt{z^2-4z-7}\right|+C$, where $a,b\in I$, then
\begin{qoptions}
\opt{A}{$a>b$}
\opt{B}{$a<b$}
\opt{C}{$a+b=117$}
\opt{D}{exactly one of $a,b$ is prime}
\end{qoptions}
\end{bankq}
\begin{bankq}[id=TOAS10-P1-MATH-S3-Q01, source={TOAS - Test 10 (12 Jun 2022)}, paper=1, qno=1, type=numerical, subtopic={Substitution}, answer={0.50}, answer_src=key, difficulty=5, diff_parts={K1 S2 I2 C0 F1}, diff_note={concepts: integration by substitution | steps: t = x\^{}2/(1 + x\^{}2) gives dt = 2x dx/(1 + x\^{}2)\^{}2; integral is t\^{}9/18; read m and n | I: choose t = x\^{}2/(1 + x\^{}2)}]
If $\displaystyle\int\dfrac{x^{17}}{(1+x^2)^{10}}\,dx=\dfrac1n\left(\dfrac{x^2}{1+x^2}\right)^m+C$, where $m,n\in N$, then find $\dfrac mn$ in lowest form.
\end{bankq}
\begin{bankq}[id=TOAS10-P1-MATH-S3-Q04, source={TOAS - Test 10 (12 Jun 2022)}, paper=1, qno=4, type=numerical, subtopic={Standard forms}, answer={4.58 to 4.59}, answer_src=key, difficulty=6, diff_parts={K2 S3 I1 C0 F1}, diff_note={concepts: trigonometric simplification, integration | steps: combine the four terms into one fraction; simplify to ((sin x + cos x - 1)/2)\^{}2; expand; integrate; read alpha, beta, gamma | I: combine into sin x cos x over 1 + sin x + cos x}]
If $\displaystyle\int\dfrac{dx}{(\sec x+\operatorname{cosec}x+\tan x+\cot x)^2}=\dfrac x\alpha+\dfrac{\cos\left(x+\frac\pi4\right)}{\beta}+\dfrac{\cos2x}{\gamma}+\delta$, then find $|\alpha+\beta+\gamma|$.
\end{bankq}
\begin{bankq}[id=TOAS11-P1-MATH-S2-Q05, source={TOAS - Test 11 (27 Jun 2022)}, paper=1, qno=5, type=multi, subtopic={Substitution}, answer={B,C}, answer_src=key, difficulty=4, diff_parts={K1 S2 I1 C0 F1}, diff_note={concepts: integral of sqrt(sec\^{}2 x + 3) | steps: rewrite with cos x; substitute sin x; integrate; fix the constants | I: substitution that rationalises}]
If $\displaystyle\int\sqrt{\sec^2x+3}\,dx=f(x)+\sqrt3\,g(x)+c$ where $c$ is constant of integration and $f(0)=\ln2$, $g\left(\dfrac\pi2\right)=\dfrac\pi3+1$, then
\begin{qoptions}
\opt{A}{$f(x)$ is an odd function}
\opt{B}{$f\left(\dfrac\pi4\right)=\ln\left(1+\sqrt5\right)$}
\opt{C}{$g\left(\dfrac{3\pi}{2}\right)+g\left(\dfrac\pi2\right)=2$}
\opt{D}{$g(x)$ is an odd function}
\end{qoptions}
\end{bankq}
\begin{bankq}[id=TOAS11-P1-MATH-S3-Q03, source={TOAS - Test 11 (27 Jun 2022)}, paper=1, qno=3, type=numerical, subtopic={Substitution}, answer={2.00}, answer_src=key, difficulty=5, diff_parts={K1 S2 I2 C0 F1}, diff_note={concepts: integral with a hidden perfect square | steps: denominator = (x\^{}3 + x\^{}2 + 1)\^{}2 + 4; numerator is the derivative of x\^{}3 + x\^{}2; standard arctan integral | I: complete the square}, also_in={TOAS - Test 12 (10 Jul 2022) P1 Q3}]
If $\displaystyle\int\dfrac{(3x^2+2x)\,dx}{x^6+2x^5+x^4+2x^3+2x^2+5}=\dfrac12\tan^{-1}(g(x))+C$, then the value of $(g(1)+g(0))$ is
\end{bankq}
\begin{bankq}[id=TOAS11-P2-MATH-S3-Q06, source={TOAS - Test 11 (27 Jun 2022)}, paper=2, qno=6, type=numerical, subtopic={By parts}, answer={1.50}, answer_src=key, difficulty=4, diff_parts={K1 S2 I0 C1 F1}, diff_note={concepts: integral of polynomial times e\^{}x | steps: assume Q e\^{}x with Q + Q' = polynomial; solve the coefficients from the top degree downward; read alpha and beta}]
If $\displaystyle\int(x^6+7x^5+6x^4+5x^3+4x^2+3x+1)e^x\,dx$ is equal to $\displaystyle\sum_{k=1}^{\alpha}\beta x^k\cdot e^x+C$ (where $C$ is constant of integration) then $\left(\dfrac{\alpha}{4\beta}\right)$ is
\end{bankq}
\begin{bankq}[id=TOAS12-P1-MATH-S1-Q01, source={TOAS - Test 12 (10 Jul 2022)}, paper=1, qno=1, type=multi, subtopic={Substitution}, answer={A,C,D}, answer_src=key, difficulty=4, diff_parts={K1 S2 I1 C0 F1}, diff_note={concepts: integration by substitution | steps: rewrite (x\^{}5 - x)\^{}(1/5)/x\^{}6 with x\^{}4 inside; substitute t = 1 - 1/x\^{}4; integrate; expand the tan polynomial | I: factor out x\^{}5}]
Identify the correct statement(s) (where $C$ is the constant of integration).
\begin{qoptions}
\opt{A}{$\displaystyle\int\frac{(x^5-x)^{1/5}}{x^6}\,dx=\frac{5}{24}\left(1-\frac{1}{x^4}\right)^{6/5}+C$}
\opt{B}{$\displaystyle\int\frac{(x^5-x)^{1/5}}{x^6}\,dx=-\frac{5}{24}\left(1-\frac{1}{x^4}\right)^{6/5}+C$}
\opt{C}{$\displaystyle\int(\tan^3x+\tan^2x+\tan x+1)\,dx=\frac{(\tan x+1)^2}{2}+C$}
\opt{D}{$\displaystyle\int(\tan^3x+\tan^2x+\tan x+1)\,dx=\frac12\sec^2x+\tan x+C$}
\end{qoptions}
\end{bankq}
\begin{bankq}[id=TOAS12-P1-MATH-S1-Q03, source={TOAS - Test 12 (10 Jul 2022)}, paper=1, qno=3, type=multi, subtopic={Substitution}, answer={A,D}, answer_src=key, difficulty=4, diff_parts={K1 S2 I1 C0 F1}, diff_note={concepts: integral of sqrt(1 + cosec x) | steps: rewrite in sin x and cos x; substitute sin x = t\^{}2 style; match the options | I: substitution that rationalises}]
$\displaystyle\int\sqrt{1+\operatorname{cosec}x}\;dx$ equals - (where $C$ is integration constant)
\begin{qoptions}
\opt{A}{$2\sin^{-1}\sqrt{\sin x}+C$}
\opt{B}{$\sqrt2\cos^{-1}\sqrt{\cos x}+C$}
\opt{C}{$C-2\sin^{-1}(1-2\sin x)$}
\opt{D}{$\cos^{-1}(1-2\sin x)+C$}
\end{qoptions}
\end{bankq}
\begin{bankq}[id=TOAS12-P1-MATH-S1-Q10, source={TOAS - Test 12 (10 Jul 2022)}, paper=1, qno=10, type=multi, subtopic={Substitution}, answer={A,C}, answer_src=key, difficulty=5, diff_parts={K1 S2 I2 C0 F1}, diff_note={concepts: integral by substitution | steps: factor x\^{}17 out of the radical; substitute t = x\^{}(-13) - 1 or similar; integrate; match p, q, r | I: pick the substitution that makes the numerator a derivative}]
If $\displaystyle\int\frac{7+6x^{13}}{\sqrt[3]{x^{17}-x^{30}}}\,dx=-\dfrac{p\,(1-x^{13})^{q/3}}{2x^{r/3}}+C$ where $p,q,r\in\mathbb{N}$ and $C$ is the constant of integration, then
\begin{qoptions}
\opt{A}{$p=3$}
\opt{B}{$p=7$}
\opt{C}{$r-q=12$}
\opt{D}{$r-q=8$}
\end{qoptions}
\end{bankq}
\begin{bankq}[id=TOAS12-P2-MATH-S1-Q06, source={TOAS - Test 12 (10 Jul 2022)}, paper=2, qno=6, type=single, subtopic={By parts}, answer={D}, answer_src=key, difficulty=4, diff_parts={K1 S2 I2 C0 F0}, diff_note={concepts: integral with x sin x + 5 cos x | steps: recognise the quotient rule form with tan x; fix the constant; evaluate f(pi) | I: integrand as a derivative of a quotient}]
Let $\displaystyle\int\dfrac{x^2+20}{(x\sin x+5\cos x)^2}\,dx=f(x)+\tan x+C$, $f(0)=0$ where $C$ denotes constant of integration. Then the value of $\pi+5f(\pi)$ is equal to
\begin{qoptions}
\opt{A}{2}
\opt{B}{4}
\opt{C}{3}
\opt{D}{0}
\end{qoptions}
\end{bankq}
\begin{bankq}[id=TOAS12-P2-MATH-S3-Q01, source={TOAS - Test 12 (10 Jul 2022)}, paper=2, qno=1, type=numerical, subtopic={By parts}, answer={1.00}, answer_src=key, difficulty=5, diff_parts={K1 S2 I2 C0 F1}, diff_note={concepts: integral of e\^{}x times a function and its derivative | steps: set f + f' equal to the algebraic factor; guess f = 1/sqrt(x\^{}2 + 1); range | I: e\^{}x (f + f') form}]
If $\displaystyle\int\dfrac{e^x\,(x^2-x+1)}{(x^2+1)\sqrt{x^2+1}}\,dx=e^x\cdot f(x)+C$, where $C$ is constant of integration, then number of integers in the range of $f(x)$ is
\end{bankq}
\begin{bankq}[id=TOAS13-P1-MATH-S3-Q07, source={TOAS - Test 13 (24 Jul 2022)}, paper=1, qno=7, type=numerical, subtopic={Substitution}, answer={2.00}, answer_src=key, difficulty=4, diff_parts={K1 S2 I1 C0 F1}, diff_note={concepts: substitution in an integral | steps: t = 2x\^{}9 + 3x\^{}6 + 6x\^{}3 with dt = 18 x\^{}2 (x\^{}6 + x\^{}3 + 1) dx; rewrite the radical as t\^{}(1/3)/x; integrate}]
If $\displaystyle\int(x^9+x^6+x^3)(2x^6+3x^3+6)^{1/3}dx=a\left(2x^9+3x^6+6x^3\right)^{4/3}+C$, then the value of $48a$ is equal to \underline{\hspace{1.5cm}}.\par
(Where $C$ is constant of integration)
\end{bankq}
\begin{bankq}[id=TOAS13-P1-MATH-S3-Q08, source={TOAS - Test 13 (24 Jul 2022)}, paper=1, qno=8, type=numerical, subtopic={Substitution}, answer={0.00}, answer_src=key, difficulty=4, diff_parts={K1 S2 I1 C0 F1}, diff_note={concepts: trigonometric integral by derivative recognition | steps: multiply through by cos\^{}2 x to get (2cos x + 1)/(cos x + 2)\^{}2; recognise the derivative of sin x/(cos x + 2); read lambda and p}]
If $\displaystyle\int\dfrac{\sec x\,(2+\sec x)}{(1+2\sec x)^2}\,dx$ is $\left(\dfrac{\cos x+2}{\lambda\sin x}\right)^p$, then $\lambda+p$ is equal to \underline{\hspace{1.5cm}}.
\end{bankq}
\begin{bankq}[id=TOAS13-P2-MATH-S1-Q09, source={TOAS - Test 13 (24 Jul 2022)}, paper=2, qno=9, type=single, subtopic={Substitution}, answer={B}, answer_src=key, difficulty=2, diff_parts={K1 S1 I1 C0 F0}, diff_note={concepts: integration by recognising a quotient derivative | steps: try sqrt(x\^{}4+2x\^{}3+1)/x; differentiate and match the numerator}]
$\displaystyle\int\dfrac{x^4+x^3-1}{x^2\sqrt{x^4+1+2x^3}}\,dx$ is equal to (where $C$ is constant of integration)
\begin{qoptions}
\opt{A}{$\dfrac{\sqrt{x^4+2x^3+1}}{x^{1/2}}+C$}
\opt{B}{$\dfrac{\sqrt{x^4+2x^3+1}}{x}+C$}
\opt{C}{$\dfrac{\sqrt{x^4+2x^3+1}}{x^{3/2}}+C$}
\opt{D}{$\dfrac{\sqrt{x^4+2x^3+1}}{x^3}+C$}
\end{qoptions}
\end{bankq}
\begin{bankq}[id=TOAS13-P2-MATH-S1-Q17, source={TOAS - Test 13 (24 Jul 2022)}, paper=2, qno=17, type=single, subtopic={Substitution}, answer={B}, answer_src=key, difficulty=2, diff_parts={K1 S1 I1 C0 F0}, diff_note={concepts: integration by recognising a derivative | steps: factor x\^{}2 under the root; the numerator is the derivative of the inside divided by 15}]
If $\displaystyle\int\dfrac{x^6(x+1)}{\sqrt{5x^{10}+6x^9+x^4}}\,dx=\dfrac{\sqrt{5x^6+6x^5+1}}{k}+C$ then $k=$
\begin{qoptions}
\opt{A}{10}
\opt{B}{15}
\opt{C}{5}
\opt{D}{None}
\end{qoptions}
\end{bankq}
\begin{bankq}[id=TOAS13-P2-MATH-S2-Q03, source={TOAS - Test 13 (24 Jul 2022)}, paper=2, qno=3, type=numerical, subtopic={Partial fractions}, answer={0.33}, answer_src=solved, working={$x^4+1=(x^2+1)^2-2x^2$ and $x^6+1=(x^2+1)(x^4-x^2+1)$, so integrand $=\dfrac{x^2+1}{x^4-x^2+1}-\dfrac{2x^2}{x^6+1}$; the first is $\dfrac d{dx}\tan^{-1}\left(x-\dfrac1x\right)$, the second gives $-\tfrac23\tan^{-1}x^3$. $A=1$, $B=-\tfrac23$, $A+B=\tfrac13$}, difficulty=5, diff_parts={K2 S2 I1 C0 F1}, diff_note={concepts: rewriting a rational integrand, standard arctan forms | steps: x\^{}4 + 1 = (x\^{}2 + 1)\^{}2 - 2x\^{}2 and factor x\^{}6 + 1; first term is d of arctan(x - 1/x); second term is arctan x\^{}3 times 2/3}]
$\displaystyle\int\dfrac{x^4+1}{x^6+1}dx=A\tan^{-1}\left(x-\dfrac1x\right)+B\tan^{-1}\left(x^3\right)+C$ then value of $(A+B)$ is
\end{bankq}
\begin{bankq}[id=TOAS13-P2-MATH-S2-Q05, source={TOAS - Test 13 (24 Jul 2022)}, paper=2, qno=5, type=numerical, subtopic={Substitution}, answer={6}, answer_src=solved, working={Put $u=1+\sqrt x$: $\int\dfrac{2(u-1)}{u^8}du=-\dfrac2{6u^6}+\dfrac2{7u^7}+C$, so $a=6$ (and $b=7$)}, difficulty=3, diff_parts={K1 S2 I0 C0 F1}, diff_note={concepts: substitution in an integral | steps: u = 1 + sqrt x so dx = 2 (u - 1) du; integrate powers of u; read a}]
If $\displaystyle\int\dfrac{1}{\left(1+\sqrt x\right)^8}dx=\dfrac{-2}{a\left(1+\sqrt x\right)^a}+\dfrac{2}{b\left(1+\sqrt x\right)^b}+C$ (where $C$ is constant of integration), then the value of `a' is equal to
\end{bankq}
\begin{bankq}[id=TOAS14-P1-MATH-S1-Q01, source={TOAS - Test 14 (07 Aug 2022)}, paper=1, qno=1, type=multi, subtopic={Substitution}, answer={C}, answer_src=key, difficulty=3, diff_parts={K1 S1 I1 C0 F1}, diff_note={concepts: integral of a trigonometric reciprocal | steps: write sin x + cos x as root 2 sin of x + pi/4; integral of 1/(1 + sin t) is tan t minus sec t}]
$\displaystyle\int\dfrac{dx}{\sqrt2+\sin x+\cos x}$ is equal to -
\begin{qoptions}
\opt{A}{$\dfrac1{\sqrt2}\left(\sec\left(x+\dfrac\pi4\right)-\tan\left(x+\dfrac\pi4\right)\right)+C$}
\opt{B}{$\dfrac1{\sqrt2}\left(\sec\left(x-\dfrac\pi4\right)-\tan\left(x-\dfrac\pi4\right)\right)+C$}
\opt{C}{$\dfrac1{\sqrt2}\left(\tan\left(x+\dfrac\pi4\right)-\sec\left(x+\dfrac\pi4\right)\right)+C$}
\opt{D}{$\dfrac1{\sqrt2}\left(\tan\left(x-\dfrac\pi4\right)-\sec\left(x-\dfrac\pi4\right)\right)+C$}
\end{qoptions}
\end{bankq}
\begin{bankq}[id=TOAS15-P2-MATH-S1-Q04, source={TOAS - Test 15 (22 Aug 2022)}, paper=2, qno=4, type=multi, subtopic={Substitution}, answer={B,C,D}, answer_src=key, difficulty=4, diff_parts={K1 S2 I1 C0 F1}, diff_note={concepts: substitution by the sixth root | steps: x = t\^{}6; reduce to a rational function; integrate; read the constants}]
If $\displaystyle\int\dfrac{x+\sqrt[3]{x^2}+\sqrt[6]{x}}{x\left(1+\sqrt[3]{x}\right)}\,dx=\alpha\cdot x^\beta+\gamma\tan^{-1}(x^\delta)+C$, where $C$ is constant of integration, then which of the following is(are) correct? $\{2\alpha,3\beta,4\gamma,6\delta\in N\}$
\begin{qoptions}
\opt{A}{$\alpha=\dfrac23$}
\opt{B}{$\alpha\beta=1$}
\opt{C}{$(\alpha\beta+\gamma)=7$}
\opt{D}{Sum of all coefficients of even powers of $x$ in the expansion of $\left(\alpha\beta+\dfrac x\delta\right)^\gamma$ is equal to $\dfrac{(\gamma+1)^\gamma+(\gamma-1)^\gamma}{2}$.}
\end{qoptions}
\end{bankq}
\begin{bankq}[id=TOAS15-P2-MATH-S1-Q06, source={TOAS - Test 15 (22 Aug 2022)}, paper=2, qno=6, type=multi, subtopic={By parts}, answer={B,C}, answer_src=key, difficulty=6, diff_parts={K2 S2 I2 C0 F1}, diff_note={concepts: derivative of an integral, integration by parts | steps: differentiate the inner integrals; the exponent integrand becomes the derivative of a product; evaluate derivatives at 0}]
If\par\noindent\hfill$F(x)=\displaystyle\int e^x\left(\int_0^xx\sin x^2dx+\int_0^x\left(\sin x^2+2x^2\cos x^2\right)dx\right)dx$\hfill\null\par\noindent$F(0)=\dfrac{-1}2$ then which of the following statement(s) must be true.
\begin{qoptions}
\opt{A}{$F\left(\sqrt{\dfrac{4\pi}7}\right)<0$}
\opt{B}{$F'(0)=\dfrac{-1}2$}
\opt{C}{$F''(0)+F'''(0)=-1$}
\opt{D}{$F(0)-F'(0)=\dfrac12$}
\end{qoptions}
\end{bankq}
\begin{bankq}[id=TOAS16-P1-MATH-S2-Q06, source={TOAS - Test 16 (04 Sep 2022)}, paper=1, qno=6, type=multi, subtopic={Standard forms}, answer={A,B,D}, answer_src=key, difficulty=2, diff_parts={K1 S1 I0 C0 F1}, diff_note={concepts: integral of a reciprocal quadratic | steps: complete the square; arctan for  | a}]
Let $\displaystyle\int\frac{dx}{x^2+ax+1}=f(g(x))+C$, where $C$ is constant of integration, then
\begin{qoptions}
\opt{A}{$f(x)$ is inverse trigonometric function for $|a|<2$}
\opt{B}{$f(x)$ is logarithmic function for $|a|>2$}
\opt{C}{$g(x)$ is quadratic function for $|a|<2$}
\opt{D}{$g(x)$ is rational function for $|a|>2$}
\end{qoptions}
\end{bankq}
\begin{bankq}[id=TOAS17-P2-MATH-S1-Q09, source={TOAS - Test 17 (18 Sep 2022)}, paper=2, qno=9, type=single, subtopic={By parts}, answer={A}, answer_src=key, difficulty=2, diff_parts={K1 S1 I1 C0 F0}, diff_note={concepts: integral of the form e\^{}x (f + f') | steps: write the integrand as e\^{}x(h + h'); evaluate with the condition}]
If $f(x)=\displaystyle\int e^x\left(\tan^{-1}(x)+\dfrac{2x}{(1+x^2)^2}\right)dx$, $f(0)=-1$ then the value of $f(1)$ is
\begin{qoptions}
\opt{A}{$e\left(\dfrac\pi4-\dfrac12\right)$}
\opt{B}{$e\left(\dfrac\pi4+\dfrac12\right)$}
\opt{C}{$e\left(\dfrac\pi2-\dfrac14\right)$}
\opt{D}{$e^{-1}\left(\dfrac\pi4-\dfrac12\right)$}
\end{qoptions}
\end{bankq}
\begin{bankq}[id=TOAS17-P2-MATH-S1-Q11, source={TOAS - Test 17 (18 Sep 2022)}, paper=2, qno=11, type=single, subtopic={By parts}, answer={A}, answer_src=key, difficulty=3, diff_parts={K1 S2 I1 C0 F0}, diff_note={concepts: reduction formula | steps: integration by parts; use the identity; recurrence}]
If $I_n=\displaystyle\int(\operatorname{cosec}2x)^n\,dx$ then
\begin{qoptions}
\opt{A}{$I_n=\dfrac{\cot2x\operatorname{cosec}^{n-2}2x}{-2(n-1)}+\dfrac{n-2}{n-1}I_{n-2}$}
\opt{B}{$I_n=\dfrac{\cot2x\operatorname{cosec}^{n-2}2x}{-2(n-1)}+\dfrac{n+2}{n-1}I_{n-2}$}
\opt{C}{$I_n=\dfrac{\cot2x\operatorname{cosec}^{n-2}2x}{-(n-1)}+\dfrac{n-2}{n-1}I_{n-2}$}
\opt{D}{$I_n=\dfrac{\cot2x\operatorname{cosec}^{n-2}2x}{-(n-1)}+\dfrac{n+2}{n-1}I_{n-2}$}
\end{qoptions}
\end{bankq}
\begin{bankq}[id=TOAS18-P2-MATH-S1-Q03, source={TOAS - Test 18 (30 Oct 2022)}, paper=2, qno=3, type=single, subtopic={Partial fractions}, answer={D}, answer_src=key, difficulty=4, diff_parts={K1 S2 I1 C1 F0}, diff_note={concepts: integration by factorisation | steps: factor numerator and denominator; partial fractions; combine logs}]
$\displaystyle\int\dfrac{x^3-x^2+x-1}{x^5+1}\,dx=f(x)$ and $f(0)=0$, then $f(x)$ is
\begin{qoptions}
\opt{A}{$\dfrac15\ell n\left(\dfrac{1+x^5}{(1+x^3)^5}\right)$}
\opt{B}{$\dfrac15\ell n\left(\dfrac{1+x^5}{1+x^3}\right)$}
\opt{C}{$\dfrac15\ell n\left(\dfrac{1+x^5}{1+x}\right)$}
\opt{D}{$\dfrac15\ell n\left(\dfrac{1+x^5}{(1+x)^5}\right)$}
\end{qoptions}
\end{bankq}
\end{document}
